% problem_id: S001-lemma-blowing-up-intersections
% source_id: S001 (Stacks Project)
% source_file: obsolete.tex
% statement_locator: obsolete.tex lines 3618--3646
% proof_locator: obsolete.tex lines 3648--3891
% env: lemma
% id_note: label 在本来源内唯一
% pairing: tight (verified)
% status: complete (statement + author proof verbatim + reason + locators)
%
\begin{lemma}
\label{lemma-blowing-up-intersections}
Let $(S, \delta)$ be as in Chow Homology, Situation \ref{chow-situation-setup}.
Let $X$ be locally of finite type over $S$.
Assume $X$ integral and $\dim_\delta(X) = d$.
Let $\{D_i\}_{i \in I}$ be a locally finite collection of
effective Cartier divisors on $X$.
Assume that for all $\{i, j, k\} \subset I$, $\#\{i, j, k\} = 3$
we have $D_i \cap D_j \cap D_k = \emptyset$.
Then there exist
\begin{enumerate}
\item an open subscheme $U \subset X$ with
$\dim_\delta(X \setminus U) \leq d - 3$,
\item a morphism $b : U' \to U$, and
\item effective Cartier divisors $\{D'_j\}_{j \in J}$ on $U'$
\end{enumerate}
with the following properties:
\begin{enumerate}
\item $b$ is proper morphism $b : U' \to U$,
\item $U'$ is integral,
\item $b$ is an isomorphism over the complement of the union of the pairwise
intersections of the $D_i|_U$,
\item $\{D'_j\}_{j \in J}$ is a locally finite collection of effective
Cartier divisors on $U'$,
\item $\dim_\delta(D'_j \cap D'_{j'}) \leq d - 2$ if $j \not = j'$, and
\item $b^{-1}(D_i|_U) = \sum n_{ij} D'_j$ for certain $n_{ij} \geq 0$.
\end{enumerate}
Moreover, if $X$ is quasi-compact, then we may assume $U = X$ in the above.
\end{lemma}

\begin{proof}
Let us first prove this in the quasi-compact case, since it is perhaps
the most interesting case. In this case we produce inductively a sequence
of blowups
$$
X = X_0 \xleftarrow{b_0} X_1 \xleftarrow{b_1} X_2 \leftarrow \ldots
$$
and finite sets of effective Cartier divisors $\{D_{n, i}\}_{i \in I_n}$.
At each stage these will have the property that any triple
intersection $D_{n, i} \cap D_{n, j} \cap D_{n, k}$ is empty.
Moreover, for each $n \geq 0$ we will have
$I_{n + 1} = I_n \amalg P(I_n)$ where $P(I_n)$ denotes
the set of pairs of elements of $I_n$. Finally, we will have
$$
b_n^{-1}(D_{n, i}) = D_{n + 1, i} +
\sum\nolimits_{i' \in I_n, i' \not = i} D_{n + 1, \{i, i'\}}
$$
We conclude that for each $n \geq 0$ we have
$(b_0 \circ \ldots \circ b_n)^{-1}(D_i)$ is a nonnegative
integer combination of the divisors $D_{n + 1, j}$, $j \in I_{n + 1}$.

\medskip\noindent
To start the induction we set $X_0 = X$ and
$I_0 = I$ and $D_{0, i} = D_i$.

\medskip\noindent
Given $(X_n, \{D_{n, i}\}_{i \in I_n})$ let $X_{n + 1}$ be the
blowup of $X_n$ in the closed subscheme
$Z_n = \bigcup_{\{i, i'\} \in P(I_n)} D_{n, i} \cap D_{n, i'}$.
Note that the closed subschemes $D_{n, i} \cap D_{n, i'}$ are pairwise
disjoint by our assumption on triple intersections.
In other words we may write
$Z_n = \coprod_{\{i, i'\} \in P(I_n)} D_{n, i} \cap D_{n, i'}$.
Moreover, in a Zariski neighbourhood of $D_{n, i} \cap D_{n, i'}$ the
morphism $b_n$ is equal to the blowup of the scheme $X_n$
in the closed subscheme $D_{n, i} \cap D_{n, i'}$, and the results
of Lemma \ref{lemma-two-divisors} apply.
Hence setting $D_{n + 1, \{i, i'\}} = b_n^{-1}(D_i \cap D_{i'})$
we get an effective Cartier divisor.
The Cartier divisors $D_{n + 1, \{i, i'\}}$ are pairwise disjoint.
Clearly we have
$b_n^{-1}(D_{n, i}) \supset D_{n + 1, \{i, i'\}}$ for
every $i' \in I_n$, $i' \not = i$. Hence, applying
Divisors, Lemma \ref{divisors-lemma-difference-effective-Cartier-divisors}
we see that indeed $b^{-1}(D_{n, i}) = D_{n + 1, i} +
\sum\nolimits_{i' \in I_n, i' \not = i} D_{n + 1, \{i, i'\}}$
for some effective Cartier divisor $D_{n + 1, i}$ on $X_{n + 1}$.
In a neighbourhood of $D_{n + 1, \{i, i'\}}$ these divisors
$D_{n + 1, i}$ play the role of the primed divisors of
Lemma \ref{lemma-two-divisors}. In particular we conclude that
$D_{n + 1, i} \cap D_{n + 1, i'} = \emptyset$ if $i \not = i'$,
$i, i' \in I_n$ by part (6) of Lemma \ref{lemma-two-divisors}.
This already implies that triple intersections
of the divisors $D_{n + 1, i}$ are zero.

\medskip\noindent
OK, and at this point we can use the quasi-compactness of $X$
to conclude that the invariant
\begin{equation}
\label{equation-invariant}
\epsilon(X, \{D_i\}_{i \in I}) =
\max\{\epsilon_Z(D_i, D_{i'}) \mid
Z \subset X,
\dim_\delta(Z) = d - 1,
\{i, i'\} \in P(I)\}
\end{equation}
is finite, since after all each $D_i$ has at most finitely many irreducible
components. We claim that for some $n$ the invariant
$\epsilon(X_n, \{D_{n, i}\}_{i \in I_n})$ is zero. Namely, if not then
by Lemma \ref{lemma-two-divisors} we have a strictly decreasing sequence
$$
\epsilon(X, \{D_i\}_{i \in I})
=
\epsilon(X_0, \{D_{0, i}\}_{i \in I_0})
>
\epsilon(X_1, \{D_{1, i}\}_{i \in I_1})
>
\ldots
$$
of positive integers which is a contradiction. Take $n$ with
invariant $\epsilon(X_n, \{D_{n, i}\}_{i \in I_n})$ equal to zero.
This means that there is no integral closed subscheme $Z \subset X_n$
and no pair of indices $i, i' \in I_n$
such that $\epsilon_Z(D_{n, i}, D_{n, i'}) > 0$.
In other words, $\dim_\delta(D_{n, i}, D_{n, i'}) \leq d - 2$ for
all pairs $\{i, i'\} \in P(I_n)$ as desired.

\medskip\noindent
Next, we come to the general case where we no longer assume that
the scheme $X$ is quasi-compact. The problem with the idea from
the first part of the proof is that we may get and infinite sequence
of blowups with centers dominating a fixed point of $X$. In order to
avoid this we cut out suitable closed subsets of codimension $\geq 3$
at each stage. Namely, we will construct by induction
a sequence of morphisms having the following shape
$$
\xymatrix{
X = X_0 \\
U_0 \ar[u]^{j_0} & X_1 \ar[l]_{b_0} \\
 & U_1 \ar[u]^{j_1} & X_2 \ar[l]_{b_1} \\
 & & U_2 \ar[u]^{j_2} & X_3 \ar[l]_{b_2}
}
$$
Each of the morphisms $j_n : U_n \to X_n$ will be an open immersion.
Each of the morphisms $b_n : X_{n + 1} \to U_n$ will be a proper birational
morphism of integral schemes. As in the quasi-compact case we will have
effective Cartier divisors $\{D_{n, i}\}_{i \in I_n}$ on $X_n$.
At each stage these will have the property that any triple
intersection $D_{n, i} \cap D_{n, j} \cap D_{n, k}$ is empty.
Moreover, for each $n \geq 0$ we will have
$I_{n + 1} = I_n \amalg P(I_n)$ where $P(I_n)$ denotes
the set of pairs of elements of $I_n$.
Finally, we will arrange it so that
$$
b_n^{-1}(D_{n, i}|_{U_n}) = D_{n + 1, i} +
\sum\nolimits_{i' \in I_n, i' \not = i} D_{n + 1, \{i, i'\}}
$$

\medskip\noindent
We start the induction by setting $X_0 = X$,
$I_0 = I$ and $D_{0, i} = D_i$.

\medskip\noindent
Given $(X_n, \{D_{n, i}\})$ we construct the open subscheme
$U_n$ as follows. For each pair $\{i, i'\} \in P(I_n)$ consider
the closed subscheme $D_{n, i} \cap D_{n, i'}$. This has ``good''
irreducible components which have $\delta$-dimension $d - 2$ and
``bad'' irreducible components which have $\delta$-dimension $d - 1$.
Let us set
$$
\text{Bad}(i, i')
=
\bigcup\nolimits_{W \subset D_{n, i} \cap D_{n, i'}
\text{ irred.\ comp. with }\dim_\delta(W) = d - 1} W
$$
and similarly
$$
\text{Good}(i, i')
=
\bigcup\nolimits_{W \subset D_{n, i} \cap D_{n, i'}
\text{ irred.\ comp. with }\dim_\delta(W) = d - 2} W.
$$
Then $D_{n, i} \cap D_{n, i'} = \text{Bad}(i, i') \cup \text{Good}(i, i')$
and moreover we have
$\dim_\delta(\text{Bad}(i, i') \cap \text{Good}(i, i')) \leq d - 3$.
Here is our choice of $U_n$:
$$
U_n
=
X_n
\setminus
\bigcup\nolimits_{\{i, i'\} \in P(I_n)}
\text{Bad}(i, i') \cap \text{Good}(i, i').
$$
By our condition on triple intersections of the divisors $D_{n, i}$
we see that the union is actually a disjoint union. Moreover,
we see that (as a scheme)
$$
D_{n, i}|_{U_n} \cap D_{n, i'}|_{U_n}
=
Z_{n, i, i'} \amalg G_{n, i, i'}
$$
where $Z_{n, i, i'}$ is $\delta$-equidimensional of dimension $d - 1$
and $G_{n, i, i'}$ is $\delta$-equidimensional of dimension $d - 2$.
(So topologically $Z_{n, i, i'}$ is the union of the bad components
but throw out intersections with good components.) Finally we set
$$
Z_n =
\bigcup\nolimits_{\{i, i'\} \in P(I_n)} Z_{n, i, i'} =
\coprod\nolimits_{\{i, i'\} \in P(I_n)} Z_{n, i, i'},
$$
and we let $b_n : X_{n + 1} \to X_n$ be the blowup in $Z_n$.
Note that Lemma \ref{lemma-two-divisors}
applies to the morphism $b_n : X_{n + 1} \to X_n$ locally around
each of the loci $D_{n, i}|_{U_n} \cap D_{n, i'}|_{U_n}$. Hence,
exactly as in the first part of the proof we obtain effective
Cartier divisors $D_{n + 1, \{i, i'\}}$ for $\{i, i'\} \in P(I_n)$
and effective Cartier divisors $D_{n + 1, i}$ for $i \in I_n$
such that
$b_n^{-1}(D_{n, i}|_{U_n}) = D_{n + 1, i} +
\sum\nolimits_{i' \in I_n, i' \not = i} D_{n + 1, \{i, i'\}}$.
For each $n$ denote $\pi_n : X_n \to X$ the morphism obtained
as the composition $j_0 \circ \ldots \circ j_{n - 1} \circ b_{n - 1}$.

\medskip\noindent
{\bf Claim:} given any quasi-compact open $V \subset X$
for all sufficiently large $n$ the maps
$$
\pi_n^{-1}(V) \leftarrow \pi_{n + 1}^{-1}(V) \leftarrow \ldots
$$
are all isomorphisms. Namely, if the map
$\pi_n^{-1}(V) \leftarrow \pi_{n + 1}^{-1}(V)$ is not an isomorphism,
then $Z_{n, i, i'} \cap \pi_n^{-1}(V) \not = \emptyset$ for some
$\{i, i'\} \in P(I_n)$. Hence there exists an irreducible component
$W \subset D_{n, i} \cap D_{n, i'}$ with $\dim_\delta(W) = d - 1$.
In particular we see that $\epsilon_W(D_{n, i}, D_{n, i'}) > 0$.
Applying Lemma \ref{lemma-two-divisors} repeatedly we see that
$$
\epsilon_W(D_{n, i}, D_{n, i'})
<
\epsilon(V, \{D_i|_V\}) - n
$$
with $\epsilon(V, \{D_i|_V\})$ as in (\ref{equation-invariant}).
Since $V$ is quasi-compact, we have
$\epsilon(V, \{D_i|_V\}) < \infty$ and taking $n > \epsilon(V, \{D_i|_V\})$
we see the result.

\medskip\noindent
Note that by construction the difference $X_n \setminus U_n$
has $\dim_\delta(X_n \setminus U_n) \leq d - 3$.
Let $T_n = \pi_n(X_n \setminus U_n)$ be its image in $X$.
Traversing in the diagram of maps above using each $b_n$ is closed
it follows that $T_0 \cup \ldots \cup T_n$ is a closed subset of $X$
for each $n$. Any $t \in T_n$ satisfies $\delta(t) \leq d - 3$
by construction. Hence $\overline{T_n} \subset X$ is a closed subset
with $\dim_\delta(T_n) \leq d - 3$. By the claim above we see
that for any quasi-compact open $V \subset X$ we have
$T_n \cap V \not = \emptyset$ for at most finitely many $n$.
Hence $\{\overline{T_n}\}_{n \geq 0}$ is a locally finite
collection of closed subsets, and we may set
$U = X \setminus \bigcup \overline{T_n}$. This will be
$U$ as in the lemma.

\medskip\noindent
Note that $U_n \cap \pi_n^{-1}(U) = \pi_n^{-1}(U)$ by construction
of $U$. Hence all the morphisms
$$
b_n : \pi_{n + 1}^{-1}(U) \longrightarrow \pi_n^{-1}(U)
$$
are proper. Moreover, by the claim they eventually become isomorphisms
over each quasi-compact open of $X$. Hence we can define
$$
U' = \lim_n \pi_n^{-1}(U).
$$
The induced morphism $b : U' \to U$ is proper since this is local
on $U$, and over each compact open the limit stabilizes. Similarly
we set $J = \bigcup_{n \geq 0} I_n$ using the inclusions
$I_n \to I_{n + 1}$ from the construction. For $j \in J$ choose
an $n_0$ such that $j$ corresponds to $i \in I_{n_0}$ and define
$D'_j = \lim_{n \geq n_0} D_{n, i}$. Again this makes sense
as locally over $X$ the morphisms stabilize.
The other claims of the lemma are verified as in the case
of a quasi-compact $X$.
\end{proof}
